\section{The layered architecture}
\label{sec:arch}

\bookfigure{architecture}{The driver stack. Every horizontal line is a boundary somebody is not
allowed to look through.}{fig:arch}

Two seams carry the whole design, and they solve different problems:

\begin{booktable}{@{}lL@{}}
\thead{Seam} & \thead{Makes what testable} \\ \midrule
\code{driver::can::Interface} & \textbf{The application.} \code{EchoNode} can be run against a
  \code{Stub} and never has to know that hardware exists. \\
\code{driver::can::ByteTransport} & \textbf{The driver itself.} The same \code{Spi} code that is
  flashed to the target can be run on your laptop against a transport you script. \\
\end{booktable}

The second seam is not obvious, and it arrives in \kapref{ch:byte}. Without it, \code{Spi} can
only be run on real hardware --- and that does not exist until \kapref{ch:bringup}, second to last
in the course, on a board somebody else is debugging at the same time.

\textbf{The rule that follows, and the one a reviewer should look for: no file above
\code{driver/can/interface.hpp} may contain the word SPI.}

\subsection{The seam has already paid off once}

The argument for an interface is usually hypothetical --- *what if the hardware is replaced*. Here
it is not hypothetical.

An earlier version of this design reached the registers through \textbf{memory-mapped I/O} at the
base address \code{0xFF200000}. That assumes a processor with an address bus out to the FPGA
fabric, and the microcontroller in this system has no such bus. The registers are reached over SPI
instead.

Everything above \code{Interface} would have survived that change \textbf{unchanged}. That is the
whole argument, demonstrated on real code.

\subsection{The transmit sequence}

Sending a frame (\code{id 0x123}, \code{dlc 2}, data \code{AA BB}) is five register writes and one
poll:

\begin{console}
1. read  STATUS        -> bit 0 set? if not: busy, return false
2. write TX_ID         = 0x00000123
3. write TX_DLC        = 0x00000002
4. write TX_DATA_LO    = 0xAABB0000
5. write TX_DATA_HI    = 0x00000000
6. write TX_SEND       = 0x00000001      -> STATUS bit 0 goes low
   ...
7. poll  STATUS        -> bit 0 back: the port is free
8. read  ERROR_FLAGS   -> 0 = the frame went out, 1 = something went wrong
\end{console}

And receiving one:

\begin{console}
1. read  STATUS        -> bit 1 set? if not: no frame, return false
2. read  RX_ID, RX_DLC
3. read  RX_DATA_LO, RX_DATA_HI, unpack, mask against RX_DLC
4. write RX_ACK        = 0x00000001      -> STATUS bit 1 goes low
\end{console}

\code{TX_SEND} \textbf{must} be written last. That write does not mean "store a value" but "start
transmitting now", and the controller then reads the other registers. Write the trigger first and
you send the previous frame's contents --- which at bring-up looks like a node running exactly one
frame behind.
